\chapter{Results: Comparing CDA and FBA}
\label{ch:results}

\begin{quote}
\small
\textbf{Status of this chapter.} This chapter is the reporting
scaffold for RQ2. Every numeric cell is currently a placeholder,
rendered as \ph, and every figure is a stub describing the plot that
will occupy that position. The section structure, the table layouts, the
metric definitions, and the reading conventions are final, so that
populating the chapter from the simulation logs is a mechanical step
rather than a rewrite. No number in this chapter should be read as a
measured result until the placeholders are replaced. \emph{This note is
to be deleted before submission.}
\end{quote}

\section{Sample and Descriptive Statistics}
\label{sec:sample}

\Cref{tab:sample} describes the replayed sample: the instruments, the
sessions, the number of messages, and the composition of the event
stream. Because both engines consume this stream unchanged, the table
characterises the input to both arms of the comparison simultaneously,
and any asymmetry in the results cannot originate here.

\begin{table}[htbp]
\centering
\caption[Sample description]{Description of the replayed sample. Counts
refer to the raw event stream, identical for both engines.}
\label{tab:sample}
\small
\begin{tabularx}{\textwidth}{@{}L c c c c@{}}
\toprule
\textbf{Characteristic} & \textbf{Liquid pair} & \textbf{Less liquid
pair} & \textbf{High-vol.\ session} & \textbf{Pooled}\\
\midrule
Sessions (days)              & \ph & \ph & \ph & \ph\\
Total messages               & \ph & \ph & \ph & \ph\\
\quad of which new orders    & \ph & \ph & \ph & \ph\\
\quad of which cancellations & \ph & \ph & \ph & \ph\\
\quad of which modifications & \ph & \ph & \ph & \ph\\
Distinct participants        & \ph & \ph & \ph & \ph\\
Trades in source data        & \ph & \ph & \ph & \ph\\
Notional volume              & \ph & \ph & \ph & \ph\\
Median inter-arrival time    & \ph & \ph & \ph & \ph\\
Realized volatility (daily)  & \ph & \ph & \ph & \ph\\
\bottomrule
\end{tabularx}
\end{table}

The median inter-arrival time in \Cref{tab:sample} deserves attention
when reading the rest of the chapter, because it determines how many
events a batch of a given length contains. If the median inter-arrival
time is of the same order as $\tau$, most batches will contain few
orders and the FBA will behave close to a delayed CDA, which is the
degenerate case anticipated in \Cref{sec:interval}.

\section{Validation Results}
\label{sec:validationresults}

\Cref{tab:validation} reports the four validation layers of
\Cref{sec:validation}. These are prerequisites, not findings: the
comparison in the following sections is only interpretable if the
control engine reproduces the venue's own output reasonably faithfully
and the batch engine converges to it in the degenerate limit.

\begin{table}[htbp]
\centering
\caption[Validation results]{Validation of the replay harness and both
engines.}
\label{tab:validation}
\small
\begin{tabularx}{\textwidth}{@{}L c c@{}}
\toprule
\textbf{Validation check} & \textbf{Result} & \textbf{Threshold}\\
\midrule
Invariant violations across all runs & \ph & 0\\
Unit tests on hand-computed clearing cases passed & \ph & all\\
CDA trade-for-trade agreement with venue output (\%) & \ph & report\\
CDA volume agreement with venue output (\%) & \ph & report\\
FBA--CDA convergence in cleared volume as $\tau \to 0$ (\% deviation) &
\ph & $\to 0$\\
Run-to-run determinism (byte-identical logs) & \ph & yes\\
\bottomrule
\end{tabularx}
\end{table}

\section{RQ2.1: Liquidity and Transaction Costs}
\label{sec:res21}

\Cref{tab:res21} reports the liquidity and cost metrics of
\Cref{tab:metrics1} for both engines, with the paired difference and its
robust standard error. Differences are computed per interval and then
aggregated, following \Cref{sec:experiments}.

\begin{table}[htbp]
\centering
\caption[RQ2.1 results]{Liquidity and transaction costs under both
mechanisms. Paired differences per interval; standard errors robust to
autocorrelation in parentheses.}
\label{tab:res21}
\small
\begin{tabularx}{\textwidth}{@{}L c c c c@{}}
\toprule
\textbf{Metric} & \textbf{CDA} & \textbf{FBA} & \textbf{Difference} &
\textbf{$p$}\\
\midrule
Quoted spread (bps)                & \ph & \ph & \ph & \ph\\
Depth at best (notional)           & \ph & \ph & \ph & \ph\\
Depth within 5 bps                 & \ph & \ph & \ph & \ph\\
Depth within 25 bps                & \ph & \ph & \ph & \ph\\
Book imbalance (abs., top 5)       & \ph & \ph & \ph & \ph\\
Effective spread (bps)             & \ph & \ph & \ph & \ph\\
Realized spread, $\Delta=1$s (bps) & \ph & \ph & \ph & \ph\\
Realized spread, $\Delta=5$s (bps) & \ph & \ph & \ph & \ph\\
Realized spread, $\Delta=30$s (bps)& \ph & \ph & \ph & \ph\\
Price impact (bps)                 & \ph & \ph & \ph & \ph\\
Amihud illiquidity ($\times 10^6$) & \ph & \ph & \ph & \ph\\
\bottomrule
\end{tabularx}
\end{table}

The interpretation hinges on the decomposition of
\Cref{sec:adverse}. The effective spread is the total cost of a trade;
the realized spread is what the liquidity provider retains after the
midpoint has adjusted; the difference is price impact, the adverse
selection component. The theory of \Cref{sec:discretetime} predicts that
if batching helps, the improvement should concentrate in price impact,
because that is the component latency arbitrage inflates. An improvement
concentrated in the realized spread instead would indicate a transfer to
or from liquidity providers rather than a reduction in adverse
selection, and would require a different explanation.

The quoted spread row must be read with the caveat established in
\Cref{sec:quality}: for the FBA it is not an observed quote but the
implied spread between the marginal unexecuted buy and the marginal
unexecuted sell of the batch. It is the closest available analogue, not
the same object.

\begin{figure}[htbp]
\centering
\fbox{\begin{minipage}{0.9\textwidth}\centering\vspace{2.0cm}
\textsf{\small [Figure stub] Effective spread decomposed into realized
spread and price impact, CDA versus FBA, stacked bars per instrument
and session.}\vspace{2.0cm}\end{minipage}}
\caption[Spread decomposition]{Decomposition of the effective spread
into the retained (realized) and adverse selection (price impact)
components under both mechanisms.}
\label{fig:spreaddecomp}
\end{figure}

\section{RQ2.2: Price Discovery and Market Quality}
\label{sec:res22}

\Cref{tab:res22} reports the price discovery metrics.

\begin{table}[htbp]
\centering
\caption[RQ2.2 results]{Price discovery and market quality under both
mechanisms.}
\label{tab:res22}
\small
\begin{tabularx}{\textwidth}{@{}L c c c c@{}}
\toprule
\textbf{Metric} & \textbf{CDA} & \textbf{FBA} & \textbf{Difference} &
\textbf{$p$}\\
\midrule
Realized volatility, 1s intervals (bps)   & \ph & \ph & \ph & \ph\\
Realized volatility, 60s intervals (bps)  & \ph & \ph & \ph & \ph\\
Intra-interval price dispersion (bps)     & \ph & 0.00 & \ph & ---\\
Pricing error vs.\ oracle, mean (bps)     & \ph & \ph & \ph & \ph\\
Pricing error vs.\ oracle, 95th pct (bps) & \ph & \ph & \ph & \ph\\
\bottomrule
\end{tabularx}
\end{table}

One row of \Cref{tab:res22} is not an empirical finding and is marked
accordingly. Intra-interval price dispersion is zero for the FBA by
construction, because every trade in an interval executes at $\clr$.
Reporting it is nonetheless worthwhile for two reasons. It is the
direct structural measure of the property that defines the difference
between the two mechanisms, so its magnitude on the CDA side quantifies
exactly how much within-interval price variation uniform clearing
eliminates. And it functions as a check on the implementation: a
non-zero value on the FBA side would indicate a bug, not a finding.
This is prediction~P1 of \Cref{tab:predictions}, and it is confirmed by
construction rather than tested.

Pricing error against the oracle is the more informative row. It asks
whether the batch clearing price tracks the external reference as
closely as the continuous midpoint does. Theory does not give an
unambiguous prediction here. Batching aggregates more orders per pricing
event, which should improve the estimate, but it also prices with a
delay of up to $\tau$, which in a fast-moving market means pricing
against information that is up to one interval old. The 95th percentile
row is included because the mean can hide exactly the fast-market
episodes where the delay matters most.

\begin{figure}[htbp]
\centering
\fbox{\begin{minipage}{0.9\textwidth}\centering\vspace{2.0cm}
\textsf{\small [Figure stub] Time series of the CDA midpoint, the FBA
clearing price, and the oracle reference price over a
high-volatility episode, with the pricing error of each mechanism in a
lower panel.}\vspace{2.0cm}\end{minipage}}
\caption[Price paths against the oracle reference]{Price paths of both
mechanisms against the external reference during a high-volatility
episode.}
\label{fig:pricepaths}
\end{figure}

\section{RQ2.3: Execution, Allocation and Order-Arrival Timing}
\label{sec:res23}

\Cref{tab:res23} reports the execution and allocation metrics. The
interpretive caveat of \Cref{sec:testable} applies throughout this
section: order size inflation and boundary concentration describe
mechanical properties of fixed historical flow, not behavioural
responses.

\begin{table}[htbp]
\centering
\caption[RQ2.3 results]{Execution, allocation and order-arrival timing
under both mechanisms, plus engine controls.}
\label{tab:res23}
\small
\begin{tabularx}{\textwidth}{@{}L c c c@{}}
\toprule
\textbf{Metric} & \textbf{CDA} & \textbf{FBA} & \textbf{Difference}\\
\midrule
Executed volume (notional)            & \ph & \ph & \ph\\
Number of trades                      & \ph & \ph & \ph\\
Fill rate, all orders (\%)            & \ph & \ph & \ph\\
Fill rate, top-decile participants (\%) & \ph & \ph & \ph\\
Median time to execution (ms)         & \ph & \ph & \ph\\
95th pct.\ time to execution (ms)     & \ph & \ph & \ph\\
Trader surplus (notional)             & \ph & \ph & \ph\\
Trader surplus per executed unit (bps)& \ph & \ph & \ph\\
Order size inflation ratio            & \ph & \ph & \ph\\
Order-to-trade ratio                  & \ph & \ph & \ph\\
Boundary concentration, final 10\% (\%) & --- & \ph & ---\\
Unexecuted residual (\% of batch volume)& --- & \ph & ---\\
\midrule
\multicolumn{4}{@{}l}{\emph{Engine controls}}\\
Throughput (orders/s)                 & \ph & \ph & \ph\\
Mean clearing / match latency ($\mu$s)& \ph & \ph & \ph\\
99th pct.\ latency ($\mu$s)           & \ph & \ph & \ph\\
\bottomrule
\end{tabularx}
\end{table}

Three rows carry most of the interpretive weight.

\textbf{Time to execution} quantifies the delay cost of batching, which
is the only unambiguous disadvantage the mechanism has. For the FBA it
is bounded below by the time from submission to the next boundary, so
its distribution is essentially the distribution of the position of
arrivals within the interval, shifted by any rollover. Comparing the
median with the 95th percentile separates the routine delay from the
tail created by orders that fail to clear and roll into later batches.

\textbf{Trader surplus} measures the price improvement relative to the
submitted limit that uniform clearing produces. Prediction~P9 says this
should be positive for every non-marginal executed order in the FBA. The
economically meaningful comparison is not the level but the difference
against the CDA, where executions occur at the resting price and the
surplus split depends on which side was resting.

\textbf{Engine controls} rule out the implementation explanation. If
throughput and latency are of the same order for both engines, a
difference in market outcomes cannot be attributed to one engine simply
being slower under load.

\section{The Interval Sweep}
\label{sec:res_sweep}

\Cref{tab:sweep} reports the headline metrics across the interval sweep
of \Cref{sec:experiments}. This is the empirical counterpart of the
trade-off argued theoretically in \Cref{sec:interval}: longer intervals
accumulate more orders per batch, which should raise cleared volume per
batch and improve the thickness of the auction, while also raising the
delay imposed on every participant and the staleness of the information
priced.

\begin{table}[htbp]
\centering
\caption[Interval sweep]{Headline metrics as a function of the batch
interval $\tau$. The CDA column is the $\tau$-independent reference.}
\label{tab:sweep}
\small
\begin{tabularx}{\textwidth}{@{}L c c c c c c c c@{}}
\toprule
\textbf{Metric} & \textbf{CDA} & \textbf{1\,ms} & \textbf{10\,ms} &
\textbf{50\,ms} & \textbf{100\,ms} & \textbf{250\,ms} &
\textbf{500\,ms} & \textbf{1\,s}\\
\midrule
Effective spread (bps)      & \ph & \ph & \ph & \ph & \ph & \ph & \ph & \ph\\
Price impact (bps)          & \ph & \ph & \ph & \ph & \ph & \ph & \ph & \ph\\
Executed volume (index)     & \ph & \ph & \ph & \ph & \ph & \ph & \ph & \ph\\
Orders per batch            & --- & \ph & \ph & \ph & \ph & \ph & \ph & \ph\\
Empty batches (\%)          & --- & \ph & \ph & \ph & \ph & \ph & \ph & \ph\\
Median time to exec.\ (ms)  & \ph & \ph & \ph & \ph & \ph & \ph & \ph & \ph\\
Pricing error (bps)         & \ph & \ph & \ph & \ph & \ph & \ph & \ph & \ph\\
Unexecuted residual (\%)    & --- & \ph & \ph & \ph & \ph & \ph & \ph & \ph\\
\bottomrule
\end{tabularx}
\end{table}

The share of empty batches is the diagnostic for the degeneracy
discussed in \Cref{sec:interval}. If a large fraction of batches contain
no crossing orders, the mechanism is not aggregating anything and the
comparison at that interval says little about batching as such.

\begin{figure}[htbp]
\centering
\fbox{\begin{minipage}{0.9\textwidth}\centering\vspace{2.0cm}
\textsf{\small [Figure stub] Two-panel plot against $\tau$ on a log
scale: upper panel, effective spread and price impact with the CDA
reference as a horizontal line; lower panel, median time to execution
and orders per batch.}\vspace{2.0cm}\end{minipage}}
\caption[The interval trade-off]{The batch interval trade-off:
transaction costs against delay, as a function of $\tau$.}
\label{fig:sweep}
\end{figure}

\section{The Rationing Rule Comparison}
\label{sec:res_rationing}

\Cref{tab:rationingres} compares the four rationing rules at the
baseline interval. Because the flow is fixed, the rules cannot differ in
clearing price or in total cleared volume; by the complementary
slackness argument of \Cref{sec:clearing}, they differ only in the
allocation of the marginal quantity at $\clr$. What the table therefore
measures is distributional: who among the marginal participants gets
filled.

\begin{table}[htbp]
\centering
\caption[Rationing rule comparison]{Allocation outcomes under the four
rationing rules, baseline $\tau$. Clearing price and total volume are
invariant to the rule by construction.}
\label{tab:rationingres}
\small
\begin{tabularx}{\textwidth}{@{}L c c c c@{}}
\toprule
\textbf{Outcome} & \textbf{Price-time} & \textbf{Pro rata} &
\textbf{Size priority} & \textbf{Random}\\
\midrule
Share of volume rationed (\%)          & \ph & \ph & \ph & \ph\\
Participants receiving partial fills   & \ph & \ph & \ph & \ph\\
Gini of marginal allocation            & \ph & \ph & \ph & \ph\\
Fill rate, largest size quartile (\%)  & \ph & \ph & \ph & \ph\\
Fill rate, smallest size quartile (\%) & \ph & \ph & \ph & \ph\\
Fill rate, earliest arrival quartile (\%) & \ph & \ph & \ph & \ph\\
Corr(fill share, submitted size)       & \ph & \ph & \ph & \ph\\
Corr(fill share, arrival position)     & \ph & \ph & \ph & \ph\\
\bottomrule
\end{tabularx}
\end{table}

The two correlation rows are the informative ones, because they measure
directly which variable the rule makes allocation depend on, and
therefore which variable participants would compete on if they could
adapt. A high correlation between fill share and submitted size under
pro rata is the mechanical precondition for the size inflation that
prediction~P5 describes; a high correlation with arrival position under
price-time priority is the mechanical precondition for the micro-scale
speed race of prediction~P6. Observing the precondition is not observing
the behaviour, and \Cref{sec:testable} explains why this design cannot
observe the latter.

\section{Summary of Findings}
\label{sec:res_summary}

\ph[to be written once results are populated] The summary will state,
for each of RQ2.1, RQ2.2 and RQ2.3, the direction and magnitude of the
measured difference, whether it is consistent with the corresponding
prediction in \Cref{tab:predictions}, and whether it is robust across
instruments, sessions and volatility regimes. It will also state
explicitly which predictions were confirmed by construction rather than
tested, so that \Cref{ch:discussion} does not overstate what the
simulation established.
